\chapter{二次型}

\section{二次型的定义}

在几何部分中，我们讨论了二次曲线和二次曲面的分类. 对于平面二次曲线
\[ax^2 + by^2 + cxy + dx + ey + f = 0,\]
常见的处理手段是通过旋转变换消去交叉项$xy$，然后平移变换消去一次项，最终得到标准形式. 类似的多项式和变换形式在数学与物理中也有很多应用，为了方便我们考虑二次齐次多项式，也就引入了二次型的概念. 需要注意的是，本讲的讨论对象是数域（主要是实数域）上的二次型；非数域的情形需要采用不同的方法，将另作讨论.

\begin{definition}{二次型}{} \index{ercixing@二次型 (quadratic form)}
    $n$个元$x_1,x_2,\ldots,x_n$的二次齐次多项式
    \begin{align*}
        f(x_1,x_2,\ldots,x_n) & = \sum_{i=1}^{n}a_{ii}x_i^2+\sum\limits_{1\leqslant i<j\leqslant n}2a_{ij}x_ix_j    \\
                              & = a_{11}x_1^2+a_{22}x_2^2+\cdots+a_{nn}x_n^2                                        \\
                              & \quad +2a_{12}x_1x_2+\cdots+2a_{1n}x_1x_n+2a_{23}x_2x_3+\cdots+2a_{n-1,n}x_{n-1}x_n
    \end{align*}
    称为域$\mathbf{F}$上的$n$元二次型（简称\term{二次型}）.
\end{definition}

若令$a_{ij}=a_{ji}\enspace(1\leqslant i<j\leqslant n)$，则二次型可表示为
\[f(x_1,x_2,\ldots,x_n)=\sum_{i=1}^{n}\sum_{j=1}^{n}a_{ij}x_ix_j=X^\mathrm{T}AX\]
其中$X=(x_1,x_2,\ldots,x_n)^\mathrm{T}\in\mathbf{R}^n$，$A=(a_{ij})_{n\times n}$为实对称矩阵，并称对称矩阵$A$为二次型$f(x_1,x_2,\ldots,x_n)$的相伴矩阵.

\begin{example}{}{}
    回答以下问题：
    \begin{enumerate}
        \item 已知$A$是一个$n$阶矩阵，则$A$为反对称矩阵的充要条件是对任意$n$元列向量$X$都有$X^\mathrm{T}AX=0$；

        \item 若二次型$f(x_1,x_2,\ldots,x_n)=X^\mathrm{T}AX$对任意的$n$元列向量$X$都有$f(x_1,x_2,\ldots,x_n)=0$，证明：$A=O$；

        \item 设二次型$f(x_1,x_2,\ldots,x_n)=X^\mathrm{T}AX,\enspace g(x_1,x_2,\ldots,x_n)=X^\mathrm{T}BX$.\\
              证明：若$f(x_1,x_2,\ldots,x_n)=g(x_1,x_2,\ldots,x_n)$，则$A=B$.
    \end{enumerate}
\end{example}

\begin{solution}
    \begin{enumerate}
        \item 先证充分性. 若 $A^\mathrm{T}=-A$，则对任意列向量 $X$ 都有
              \[
                  (X^\mathrm{T}AX)^\mathrm{T}=X^\mathrm{T}A^\mathrm{T}X=-X^\mathrm{T}AX.
              \]
              由于 $X^\mathrm{T}AX$ 是一个 $1$ 阶矩阵，也就是一个数，它等于自身的转置，于是
              \[
                  X^\mathrm{T}AX=-X^\mathrm{T}AX,
              \]
              从而 $X^\mathrm{T}AX=0$.

              再证必要性. 若对任意 $X$ 都有 $X^\mathrm{T}AX=0$，取 $X=e_i$ 得 $a_{ii}=0$. 再取 $X=e_i+e_j\enspace(i\neq j)$，则
              \[
                  0=(e_i+e_j)^\mathrm{T}A(e_i+e_j)=a_{ii}+a_{ij}+a_{ji}+a_{jj}=a_{ij}+a_{ji}.
              \]
              故 $a_{ji}=-a_{ij}$ 对任意 $i,j$ 都成立，即 $A^\mathrm{T}=-A$.

        \item 因为 $f(X)=X^\mathrm{T}AX$ 是二次型，所以这里的 $A$ 是对称矩阵. 已知对任意 $X$ 都有 $X^\mathrm{T}AX=0$，取 $X=e_i$，得到 $a_{ii}=0$. 再取 $X=e_i+e_j\enspace(i\neq j)$，则
              \[
                  0=(e_i+e_j)^\mathrm{T}A(e_i+e_j)=a_{ii}+2a_{ij}+a_{jj}=2a_{ij}.
              \]
              故 $a_{ij}=0$. 因而 $A=O$.

        \item 由 $f(X)=g(X)$ 可知对任意 $X$ 都有
              \[
                  X^\mathrm{T}(A-B)X=0.
              \]
              又 $A,B$ 都是二次型的相伴矩阵，所以它们都是对称矩阵，从而 $A-B$ 也是对称矩阵. 由第 2 问立刻得到 $A-B=O$，即 $A=B$.
    \end{enumerate}
\end{solution}

\begin{example}{}{}
    已知二次型
    \[f(X)=(x_1,x_2,x_3,x_4)\begin{pmatrix}
            1 & 2 & 3 & -4 \\ 3 & 2 & 1 & 4 \\ -4 & 3 & -7 & 2 \\ 0 & -6 & 8 & 4
        \end{pmatrix}\begin{pmatrix}
            x_1 \\ x_2 \\ x_3 \\ x_4
        \end{pmatrix}\]
    写出二次型$f(X)$的相伴矩阵.
\end{example}

\begin{solution}
    设题中给出的矩阵为
    \[
        M=\begin{pmatrix}
            1 & 2 & 3 & -4 \\
            3 & 2 & 1 & 4 \\
            -4 & 3 & -7 & 2 \\
            0 & -6 & 8 & 4
        \end{pmatrix}.
    \]
    因为
    \[
        X^\mathrm{T}MX=X^\mathrm{T}\left(\frac{M+M^\mathrm{T}}{2}\right)X+X^\mathrm{T}\left(\frac{M-M^\mathrm{T}}{2}\right)X,
    \]
    而 $ \dfrac{M-M^\mathrm{T}}{2} $ 是反对称矩阵，故对任意 $X$ 都有
    \[
        X^\mathrm{T}\left(\frac{M-M^\mathrm{T}}{2}\right)X=0.
    \]
    因此二次型 $f(X)$ 的相伴矩阵就是 $M$ 的对称部分
    \[
        A=\frac{M+M^\mathrm{T}}{2}
        =\begin{pmatrix}
            1 & \dfrac{5}{2} & -\dfrac{1}{2} & -2 \\
            \dfrac{5}{2} & 2 & 2 & -1 \\
            -\dfrac{1}{2} & 2 & -7 & 5 \\
            -2 & -1 & 5 & 4
        \end{pmatrix}.
    \]
\end{solution}

实际上一个形如$f(x_1,x_2,\ldots,x_n)=\displaystyle\sum_{i=1}^{n}\sum_{j=1}^{n}a_{ij}x_ix_j$的函数可以对应的矩阵是很多的，但不唯一的矩阵表示会给我们带来很多不便. 而选取对称矩阵作为二次型的相伴矩阵，不仅是因为简化了计算，还是因为二次型和对称矩阵之间存在一一对应的关系.

\section{矩阵相合的定义与性质}

二次型的基本问题是寻找一个线性变换将其变换为只包含平方项的标准形，这反映在矩阵上便是对角化. 而我们将二次曲线标准化的操作抽象出来，便会发现这是我们熟悉的坐标变换. 考虑原坐标为$X$，新坐标为$Y$，原二次型为$f(X) = X^\mathrm{T}AX$. 因为变换是线性的，并且坐标长度不变，因此我们有$X=PY$，其中$P$是可逆矩阵. 而用$Y$代换$X$，便可以得到$g(Y) = Y^\mathrm{T}(P^\mathrm{T}AP)Y$，虽然未必能立刻得到标准形，但这一变换是二次型的一个基本操作，我们称之为相合变换.

\begin{definition}{}{}
    我们称$n$阶矩阵$A$相合于$B$（记作$A\simeq B$），如果存在可逆矩阵$C$使得$B=C^\mathrm{T}AC$.
\end{definition}

矩阵相合（合同）有如下基本性质，证明我们留作习题.

\begin{enumerate}
    \item 合同是等价关系；

    \item 合同与域有关，即 $A \simeq B$ 可能在一个域上成立，在另一个域上不成立；

    \item $A \simeq B$ 一般不能得到 $A^m \simeq B^m$，但是 $A, B$ 为实对称矩阵时可以；

    \item 如果 $A, B$ 可逆且 $A \simeq B$，有 $A^{-1} \simeq B^{-1}$；

    \item 如果 $A_1\simeq A_2,B_1\simeq B_2$，则有 $\begin{pmatrix}
                  A_1 & O \\ O & B_1
              \end{pmatrix}\simeq\begin{pmatrix}
                  A_2 & O \\ O & B_2
              \end{pmatrix}$；

    \item $A \simeq B$ 表明 $A$ 可以每次做相同的初等行列变换得到 $B$，反之亦然（下面的\autoref{lem:初等变换与合同变换引理1}会给出一部分证明，另一部分留作习题）.
\end{enumerate}

\section{二次型标准形的定义与求解}

求解二次型的标准形就是寻找可逆矩阵 $C$ 使得 $C^\mathrm{T}AC$ 为对角矩阵. 最基础的想法就是如同相抵标准型那样，通过初等变换将矩阵变换为对角矩阵.

\begin{lemma}{}{初等变换与合同变换引理1}
    对称矩阵 $A$ 的下列变换都是相合变换：
    \begin{enumerate}
        \item 对换 $A$ 的第 $i$ 行与第 $j$ 行，再对换第 $i$ 列与第 $j$ 列；

        \item 将非零数 $c$ 乘以 $A$ 的第 $i$ 行，再将 $c$ 乘以 $A$ 的第 $i$ 列；

        \item 将 $A$ 的第 $i$ 行乘以 $c$ 加到第 $j$ 行，再将 $A$ 的第 $i$ 列乘以 $c$ 加到第 $j$ 列.
    \end{enumerate}
\end{lemma}

\begin{proof}
    \begin{enumerate}
        \item 设 $E_{ij}$ 为初等对换矩阵（见\autoref{def:初等矩阵}），则 $E_{ij}$ 可逆且 $E_{ij}^\mathrm{T} = E_{ij}$，而 $E_{ij}^\mathrm{T}AE_{ij}$ 恰好对应于同时交换 $A$ 的第 $i,j$ 行与第 $i,j$ 列.

        \item 设 $E_i(c)$ 为初等倍乘矩阵，则 $E_i(c)$ 可逆且 $E_i(c)^\mathrm{T}=E_i(c)$. 直接计算可知 $E_i(c)^\mathrm{T}AE_i(c)$ 恰好对应于将 $A$ 的第 $i$ 行乘以 $c$，再将第 $i$ 列乘以 $c$.

        \item 设 $E_{ij}(c)$ 为初等倍加矩阵，则 $E_{ij}(c)$ 可逆且 $E_{ij}(c)^\mathrm{T}=E_{ji}(c)$. 直接计算可知 $E_{ij}(c)^\mathrm{T}AE_{ij}(c)$ 恰好对应于将 $A$ 的第 $i$ 行乘以 $c$ 加到第 $j$ 行，再将第 $i$ 列乘以 $c$ 加到第 $j$ 列.
    \end{enumerate}
    因而上述三类变换都是合同变换.
\end{proof}

接下来的问题是，如何进行合适的初等变换来将一个对称矩阵对角化.

\begin{lemma}{}{初等变换与合同变换引理2}
    设 $A$ 是非零对称矩阵，则必定存在可逆矩阵 $C$ 使得 $C^\mathrm{T}AC$ 的第 $(1, 1)$ 个元素不为 $0$.
\end{lemma}

\begin{proof}
    \begin{enumerate}
        \item 若 $a_{11} \neq 0$，则取 $C$ 为单位矩阵，结论得证.
        \item 若 $a_{11} = 0$ 而存在 $i$ 使得 $a_{ii} \neq 0$，则将 $A$ 的第 $1$ 行与第 $i$ 行对换，再将第 $1$ 列与第 $i$ 列对换，即可使得 $a_{11} \neq 0$.
        \item 若 $a_{ii} = 0, \forall i$，那么设 $a_{ij} \neq 0(i \neq j)$，将 $A$ 的第 $j$ 行加到第 $i$ 行，再将第 $j$ 列加到第 $i$ 列. 因为 $A$ 是对称矩阵，所以 $a_{ij} = a_{ji} \neq 0$，第 $(i, i)$ 元素在变换后变为 $2 a_{ij} \neq 0$，便可根据第一种情况处理.
    \end{enumerate}
    而根据\autoref{lem:初等变换与合同变换引理1}，这些变换都是相合变换，结论得证.
\end{proof}

由此，便可以得到以下结论.

\begin{theorem}{}{}
    设 $A$ 是 $n$ 阶对称矩阵，则必存在 $\mathbf{F}$ 上的 $n$ 阶可逆矩阵 $C$，使得 $C^\mathrm{T}AC$ 为对角矩阵.
\end{theorem}

\begin{proof}
    由\autoref{lem:初等变换与合同变换引理2}，不妨设 $A = (a_{ij})$ 中 $a_{11} \neq 0$，若 $a_{i1} \neq 0$，则可将第 $1$ 行乘以 $-a_{11}^{-1}a_{i1}$ 加到第 $i$ 行上，再将第 $1$ 列乘以 $-a_{11}^{-1}a_{i1}$ 加到第 $i$ 列上，因为 $a_{i1} = a_{1i}$，所以新得到的矩阵的第 $(i, 1)$ 元素和第 $(1, i)$ 元素均等于 $0$，并且和 $A$ 相合. 重复这一过程，便可以将 $A$ 的第 $1$ 列和第 $1$ 行的非对角元素全部变为 $0$，所以 $A$ 相合与以下形式的矩阵：
    \[
        \begin{pmatrix}
            a_{11} & 0 & 0 & \cdots & 0 \\
            0 & b_{22} & b_{23} & \cdots & b_{2n} \\
            0 & b_{32} & b_{33} & \cdots & b_{3n} \\
            \vdots & \vdots & \vdots & \ddots & \vdots \\
            0 & b_{n2} & b_{n3} & \cdots & b_{nn}
        \end{pmatrix}.
    \]
    右下角是一个 $n - 1$ 阶对称矩阵，记为 $A_1$，根据归纳假设，存在 $n - 1$ 阶可逆矩阵 $D$ 使得 $D^\mathrm{T}A_1D$ 为对角矩阵，所以
    \[
        \begin{pmatrix}
            1 & O \\
            O & D^\mathrm{T}
        \end{pmatrix}
        \begin{pmatrix}
            a_{11} & O \\
            O & A_1
        \end{pmatrix}
        \begin{pmatrix}
            1 & O \\
            O & D
        \end{pmatrix} = \begin{pmatrix}
            a_{11} & O \\
            O & D^\mathrm{T}A_1D
        \end{pmatrix}
    \]

    是一个对角阵，并且

    \[
        \begin{pmatrix}
            1 & O \\
            O & D
        \end{pmatrix}^\mathrm{T} = \begin{pmatrix}
            1 & O \\
            O & D^\mathrm{T}
        \end{pmatrix}
    \]

    所以 $A$ 相合于对角阵.
\end{proof}

注意到整个证明过程暗含了如何利用初等矩阵得到相合标准形的操作，但如果想要得到坐标变换矩阵 $C$ 还需要一点额外的操作.

\begin{example}{}{相合标准形例1}
    利用初等变换法将二次型 $f(x_1,x_2,x_3)=2x_1x_2-2x_1x_3+2x_2x_3$ 化为标准形，并且求出坐标变换矩阵$C$.
\end{example}

\begin{solution}
    类似于利用初等变换求矩阵的逆，我们将相伴矩阵和单位矩阵拼接为 $\begin{pmatrix}
        A \\ E
    \end{pmatrix}$，对 $A$ 进行初等行列变换，使其变为对角矩阵，同时对 $E$ 进行对应的列变换，最终得到的 $\begin{pmatrix}
        \Lambda \\ C
    \end{pmatrix}$ 就是我们要求的结果.

    \begin{align*}
        \left(\begin{array}{c}
            A \\ \hdashline E
        \end{array}\right) ={} & \left(\begin{array}{ccc}
            0 & 1 & -1 \\
            1 & 0 & 1  \\
            -1 & 1 & 0 \\
            \hdashline
            1 & 0 & 0 \\
            0 & 1 & 0 \\
            0 & 0 & 1
        \end{array}\right)
        \xrightarrow[r_1+r_2]{c_1+c_2} \left(\begin{array}{ccc}
            2 & 1 & 0 \\
            1 & 0 & 1  \\
            0 & 1 & 0 \\
            \hdashline
            1 & 0 & 0 \\
            1 & 1 & 0 \\
            0 & 0 & 1
        \end{array}\right) & \\
        \xrightarrow[r_2-\dfrac{1}{2}r_1]{c_2-\dfrac{1}{2}c_1} & \left(\begin{array}{ccc}
            2 & 0 & 0 \\
            0 & -\dfrac{1}{2} & 1  \\
            0 & 1 & 0 \\
            \hdashline
            1 & -\dfrac{1}{2} & 0 \\
            1 & \dfrac{1}{2} & 0 \\
            0 & 0 & 1
        \end{array}\right)
        \xrightarrow[r_3+2r_2]{c_3+2c_2} \left(\begin{array}{ccc}
            2 & 0 & 0 \\
            0 & -\dfrac{1}{2} & 0  \\
            0 & 0 & 2 \\
            \hdashline
            1 & -\dfrac{1}{2} & -1 \\
            1 & \dfrac{1}{2} & 1 \\
            0 & 0 & 1
        \end{array}\right) & \\
    \end{align*}
\end{solution}

初等变换法虽然泛用性比较广，但是操作起来比较繁琐，通常情况下使用配方法更加方便. 配方法的思想非常简单，就是利用配方消除混合乘积项，将二次型表示成几个平方和的形式，最后通过坐标变换$X=CY$化标准形，其中$C$是可逆矩阵.

\begin{example}{}{相合标准形例2}
    用配方法把三元二次型
    \[f(x_1,x_2,x_3)=2x_1^2+3x_2^2+x_3^2+4x_1x_2-4x_1x_3-8x_2x_3\]
    化为标准形，并求所用的坐标变换$X=CY$即变换矩阵$C$.
\end{example}

\begin{solution}
    直接配方可得
    \begin{align*}
        f(x_1,x_2,x_3)
         & = 2x_1^2+4x_1x_2-4x_1x_3+3x_2^2-8x_2x_3+x_3^2 \\
         & = 2(x_1+x_2-x_3)^2+x_2^2-4x_2x_3-3x_3^2         \\
         & = 2(x_1+x_2-x_3)^2+(x_2-2x_3)^2-5x_3^2.
    \end{align*}
    令
    \[
        \begin{cases}
            y_1 = x_1+x_2-x_3, \\
            y_2 = x_2-2x_3,    \\
            y_3 = x_3,
        \end{cases}
    \]
    则
    \[
        f(y_1,y_2,y_3)=2y_1^2+y_2^2-5y_3^2.
    \]
    由上式反解得
    \[
        \begin{cases}
            x_3 = y_3,             \\
            x_2 = y_2+2y_3,        \\
            x_1 = y_1-y_2-y_3,
        \end{cases}
    \]
    因而所求坐标变换 $X=CY$ 为
    \[
        \begin{pmatrix}
            x_1 \\
            x_2 \\
            x_3
        \end{pmatrix}
        =
        \begin{pmatrix}
            1 & -1 & -1 \\
            0 & 1 & 2   \\
            0 & 0 & 1
        \end{pmatrix}
        \begin{pmatrix}
            y_1 \\
            y_2 \\
            y_3
        \end{pmatrix}.
    \]
\end{solution}

配方法的合理性是可逆的坐标变换所保证的，只要配方成功的消除了所有交叉项，便可以得到标准形. 但配方法在处理一些情况的时候需要一些技巧，比如\autoref{ex:相合标准形例1}，需要先通过坐标变换构造出平方项，再通过配方法消除交叉项.

\begin{example}{}{相合标准形例3}
    用配方法把二次型
    \[f(x_1,x_2,x_3)=2x_1x_2-2x_1x_3+2x_2x_3\]
    化为标准形，并求所用的坐标变换$X=CY$即变换矩阵$C$.
\end{example}

\begin{solution}
    对于 $f(x_1,x_2,x_3)=2x_1x_2-2x_1x_3+2x_2x_3$，原式中没有平方项，因此不能直接像通常那样配方. 这时先作一个可逆线性变换
    \[
        \begin{cases}
            y_1 = x_1+x_2, \\
            y_2 = x_1-x_2, \\
            y_3 = x_3.
        \end{cases}
    \]
    由此反解得
    \[
        \begin{cases}
            x_1 = \dfrac{y_1+y_2}{2}, \\
            x_2 = \dfrac{y_1-y_2}{2}, \\
            x_3 = y_3.
        \end{cases}
    \]
    代入原二次型，有
    \[ f(y_1, y_2, y_3) = \frac{1}{2}y_1^2-\frac{1}{2}y_2^2-2y_2y_3. \]
    经过这一步，二次型中已经出现了平方项，于是便可以继续配方：
    \[ f(y_1, y_2, y_3) = \frac{1}{2}y_1^2-\frac{1}{2}(y_2+2y_3)^2+2y_3^2. \]
    再令
    \[
        z_1=y_1,\qquad z_2=y_2+2y_3,\qquad z_3=y_3,
    \]
    便得到标准形
    \[
        f(z_1,z_2,z_3)=\frac{1}{2}z_1^2-\frac{1}{2}z_2^2+2z_3^2.
    \]
    如果写成从新坐标到旧坐标的形式，则由
    \[
        y_1=z_1,\qquad y_2=z_2-2z_3,\qquad y_3=z_3
    \]
    以及
    \[
        x_1=\frac{y_1+y_2}{2},\qquad x_2=\frac{y_1-y_2}{2},\qquad x_3=y_3
    \]
    可得
    \[
        \begin{cases}
            x_1 = \dfrac{1}{2}z_1+\dfrac{1}{2}z_2-z_3, \\
            x_2 = \dfrac{1}{2}z_1-\dfrac{1}{2}z_2+z_3, \\
            x_3 = z_3.
        \end{cases}
    \]
    因此坐标变换矩阵可取为
    \[
        X=
        \begin{pmatrix}
            \dfrac{1}{2} & \dfrac{1}{2} & -1 \\
            \dfrac{1}{2} & -\dfrac{1}{2} & 1 \\
            0 & 0 & 1
        \end{pmatrix}
        Z.
    \]
    这说明：当原二次型没有平方项、无法直接配方时，可以先通过一个简单的可逆线性变换人为造出平方项，再继续做配方法.
\end{solution}

不过需要注意的是，以上两种方法得到的对角矩阵不一定是相伴矩阵的相似对角化，相似对角化要求是 $P^{-1}AP$ 为对角矩阵，而相合对角化是 $P^\mathrm{T}AP$ 为对角矩阵，二者一致的话需要满足 $P^{-1} = P^\mathrm{T}$，也就是说 $P$ 是正交矩阵. 不过这一点并不是难以做到的，因为 $A$ 是实对称矩阵，根据\nameref{thm:实谱定理}，$A$ 的标准化后的特征向量组成的矩阵便是正交矩阵$P$，所以我们也可以通过正交变换求解二次型的相合标准形.

\begin{example}{}{正交变换法例1}
    用正交变换法将二次型 $f(x_1,x_2) = 2x_1^2 + 5x_2^2 + 4x_1x_2$ 化为标准形，并给出变换矩阵 $P$.
\end{example}

\begin{solution}
    二次型的相伴矩阵为
    \[
        A=\begin{pmatrix}
            2 & 2 \\
            2 & 5
        \end{pmatrix}.
    \]
    其特征多项式为
    \[
        |\lambda E-A|=\begin{vmatrix}
            \lambda-2 & -2        \\
            -2        & \lambda-5
        \end{vmatrix}
        =\lambda^2-7\lambda+6=(\lambda-1)(\lambda-6),
    \]
    故特征值为 $1,6$.

    对于特征值 $1$，解方程 $(A-E)X=0$ 得特征向量可取为 $(-2,1)^\mathrm{T}$；对于特征值 $6$，解方程 $(A-6E)X=0$ 得特征向量可取为 $(1,2)^\mathrm{T}$. 将它们单位化，得到两个单位特征向量
    \[
        u_1=\frac{1}{\sqrt{5}}\begin{pmatrix}
            -2 \\
            1
        \end{pmatrix},
        \qquad
        u_2=\frac{1}{\sqrt{5}}\begin{pmatrix}
            1 \\
            2
        \end{pmatrix}.
    \]
    令
    \[
        P=\begin{pmatrix}
            -\dfrac{2}{\sqrt{5}} & \dfrac{1}{\sqrt{5}} \\
            \dfrac{1}{\sqrt{5}}  & \dfrac{2}{\sqrt{5}}
        \end{pmatrix},
    \]
    则 $P$ 为正交矩阵，且
    \[
        P^\mathrm{T}AP=\diag(1,6).
    \]
    因此在正交变换 $X=PY$ 下，
    \[
        f(X)=Y^\mathrm{T}(P^\mathrm{T}AP)Y=y_1^2+6y_2^2.
    \]
    这就是所求标准形.
\end{solution}

\section{相合规范形 \quad 惯性定理}

从前面求解相合标准形的例子可以发现，使用配方法、初等变换法得到的相合标准形，会因为配方方式或者初等变换顺序不同而各不相同，所以相合标准形并不是唯一的. 但相抵标准形是唯一的，相似标准形不考虑特征值排列组合因素也是唯一的，因此我们也希望统一相合标准形.

不难发现，任一对角矩阵一定相合于$\diag(1,\ldots,1,-1,\ldots,-1,0,\ldots,0)$，我们称这一相合标准形为相合规范形，其中$+1$的个数称为矩阵的正惯性指数，$-1$的个数称为矩阵的负惯性指数. 并且由于变换矩阵可逆，根据相抵标准形的结论，我们有原矩阵$A$的秩$r(A)$等于这一对角矩阵的秩，于是也等于正负惯性指数之和. 显然，$A$可逆时，其相合规范形主对角元没有0.

但我们没有说明一个矩阵的相合规范形是否唯一，实际上这就是下面惯性定理的结果：
\begin{theorem}{惯性定理}{}
    实对称矩阵的相合规范形唯一.
\end{theorem}
这一定理有很多等价表述，例如实对称矩阵正、负惯性指数唯一，或者实对称矩阵相合标准形中对角线上正、负、零的个数唯一. 或者实对称矩阵特征值中正、负、零的个数唯一等.

\begin{proof}
    采取证明等价表述：实对称矩阵正、负惯性指数唯一.

    设 $f(X) = X^\mathrm{T}AX$ 为一实二次型，通过可逆坐标变换$X = BY$ 和 $X = CZ$ 得到了相合规范形
    \begin{align*}
        f ={} & y_1^2 + \cdots + y_p^2 - y_{p+1}^2 - \cdots - y_r^2 \\
          ={} & z_1^2 + \cdots + z_k^2 - z_{q+1}^2 - \cdots - z_r^2.
    \end{align*}
    其中 $r$ 为 $A$ 的秩，我们希望证明此时一定有$p = k$，使用反证法.

    假设 $p > k$，由 $Z = C^{-1}X = C^{-1}BY$，定义 $D = C^{-1}B$，则 $Z = DY$，写作线性方程组的形式为

    \[
        \begin{cases}
            z_1 & = d_{11}y_1 + \cdots + d_{1n}y_n, \\
            z_2 & = d_{21}y_1 + \cdots + d_{2n}y_n, \\
            & \vdotswithin{=} \\
            z_k & = d_{k1}y_1 + \cdots + d_{kn}y_n, \\
            & \vdotswithin{=} \\
            z_{n} & = d_{n1}y_1 + \cdots + d_{nn}y_n.
        \end{cases}
    \]

    矛盾的构造点在于 $f$ 的取值，比较明显的构造便是一种表达下为正，而另一种表达下为负. 所以令 $z_1 = z_2 = \cdots = z_k = 0$，以及 $y_{p + 1} = y_{p + 2} = \cdots = y_n = 0$. 结合上面列出的方程组，我们可以得到新的方程组.

    \[
        \begin{cases}
            d_{11}y_1 + d_{12}y_2 + \cdots + d_{1n}y_n & = 0, \\
            d_{21}y_1 + d_{22}y_2 + \cdots + d_{2n}y_n & = 0, \\
            & \vdotswithin{=} \\
            d_{k1}y_1 + d_{k2}y_2 + \cdots + d_{kn}y_n & = 0. \\
            y_{p + 1} = 0, \\
            y_{p + 2} = 0, \\
            \cdots \\
            y_n = 0.
        \end{cases}
    \]

    这一方程组有 $n$ 个变量，但只有 $k + (n - p) = n - (p - k) < n$ 个方程，所以必定有非零解，而 $y_{p + 1} = y_{p + 2} = \cdots = y_n = 0$，说明 $y_1, y_2, \ldots, y_p$ 不全为 $0$. 代入 $f$ 的表达式，可以得到

    \[
        f = y_1^2 + y_2^2 + \cdots + y_p^2 > 0.
    \]

    将 $Y$ 代入 $Z$ 关于 $Y$ 的方程中，再将得到的 $Z$ 代入 $f$ 的表达式，因为 $z_1 = z_2 = \cdots = z_k = 0$，所以

    \[
        f = -z_{k + 1}^2 - z_{k + 2}^2 - \cdots - z_n^2 \leqslant 0.
    \]

    这两者显然矛盾，所以 $p > k$ 不成立，必然有 $p \leqslant k$，同理可证 $k \leqslant p$，所以 $p = k$，也就有实对称矩阵的相合规范形唯一.

\end{proof}

\begin{example}{}{}
    确定二次型$f(x_1,x_2,\ldots,x_{10})=x_1x_2+x_3x_4+x_5x_6+x_7x_8+x_9x_{10}$的秩以及正、负惯性指数.
\end{example}

\begin{solution}
    对每一对变量分别配方：
    \[
        x_{2k-1}x_{2k}=\frac{1}{4}\bigl((x_{2k-1}+x_{2k})^2-(x_{2k-1}-x_{2k})^2\bigr), \quad k=1,2,\ldots,5.
    \]
    因而
    \[
        f=\frac{1}{4}\sum_{k=1}^5\bigl((x_{2k-1}+x_{2k})^2-(x_{2k-1}-x_{2k})^2\bigr).
    \]
    令
    \[
        y_{2k-1}=x_{2k-1}+x_{2k}, \qquad y_{2k}=x_{2k-1}-x_{2k}, \qquad k=1,2,\ldots,5,
    \]
    则这是一个可逆线性变换，并且
    \[
        f=\frac{1}{4}(y_1^2-y_2^2+y_3^2-y_4^2+y_5^2-y_6^2+y_7^2-y_8^2+y_9^2-y_{10}^2).
    \]
    因而该二次型的相合规范形中有 $5$ 个正项和 $5$ 个负项，没有零项，所以其秩为 $10$，正惯性指数为 $5$，负惯性指数为 $5$.
\end{solution}

惯性定理的``惯性''二字与物理中的惯性有关，实际上透露着某种不变性. 根据惯性定理，我们有如下结论：
\begin{enumerate}
    \item 我们可以按相合关系对全体$n$阶实对称矩阵分类，因为实对称矩阵相合意味着规范形唯一，我们可以按照$+1$、$-1$、0个数的不同划分为$\vphantom{\cfrac{n+1}{2}}\dfrac{(n+1)(n+2)}{2}$个等价类；

    \item 实域上两个实对称矩阵相合的充要条件是它们有相同的正负惯性指数，两个对角矩阵相合的充要条件是对角线上正、负、零个数相同.
\end{enumerate}

而对于复二次型，若其已经化为标准形
\[
    f(X) = d_1x_1^2 + d_2x_2^2 + \cdots + d_rx_r^2,
\]
其中 $d_i \neq 0$，$r$ 为该矩阵的秩. 设 $d_j = r_j(\cos \theta_j + \i \sin \theta_j)$，$0 \leqslant \theta_j < 2 \pi$ 那么根据棣莫弗公式（De Moivre's Formula），有
\[
    \left(\pm \sqrt{r_j}\left(\cos \dfrac{\theta_j}{2} + \i \sin \dfrac{\theta_j}{2}\right)\right)^2 = d_j.
\]
记 $\sqrt{r_j}\left(\cos \dfrac{\theta_j}{2} + \i \sin \dfrac{\theta_j}{2}\right)$ 为 $\sqrt{d_j}$，那么做如下的坐标变换$X=CY$，其中
\[C = \diag\left(\dfrac{1}{\sqrt{d_1}}, \ldots, \dfrac{1}{\sqrt{d_r}}, 1, \ldots, 1\right)\]
便可以将复二次型化为标准形
\[
    f(Y) = y_1^2 + \cdots + y_r^2.
\]
所以复二次型的标准形是唯一的，只与其秩有关.

\section{正定二次型和正定矩阵}

有一类特殊的二次型是我们非常关心的，那就是正定二次型. 这一类二次型在很多领域都有重要的应用，例如在优化问题中，正定二次型常常用来描述目标函数的性质，在统计学中，协方差矩阵必须是正定矩阵等等. 因此我们有必要对正定二次型和正定矩阵进行专门的研究.

\begin{definition}{}{}
    设 $f(X) = X^\mathrm{T}AX$ 是 $n$ 元实二次型，$A$ 是相伴矩阵.
    \begin{enumerate}
        \item 若对任意非零 $n$ 元实列向量 $X$，都有 $f(X) > 0$，则称 $f(X)$ 是正定二次型，$A$ 是正定矩阵；
        \item 若对任意非零 $n$ 元实列向量 $X$，都有 $f(X) \geqslant 0$，则称 $f(X)$ 是半正定二次型，$A$ 是半正定矩阵；
        \item 若对任意非零 $n$ 元实列向量 $X$，都有 $f(X) < 0$，则称 $f(X)$ 是负定二次型，$A$ 是负定矩阵；
        \item 若对任意非零 $n$ 元实列向量 $X$，都有 $f(X) \leqslant 0$，则称 $f(X)$ 是半负定二次型，$A$ 是半负定矩阵；
        \item 若存在 $n$ 元实列向量 $X_1, X_2$，使得 $f(X_1) > 0$ 且 $f(X_2) < 0$，则称 $f(X)$ 是不定二次型.
    \end{enumerate}
\end{definition}

其中我们重点研究正定和半正定二次型，因为负定和半负定二次型可以通过符号变换得到.

很自然的一个问题是如何判断一个二次型是正定或半正定的，我们也期望使用二次型内禀的性质来判断矩阵的性质. 显然
\[f(X) = x_1^2 + x_2^2 + \cdots + x_n^2\]
是一个正定二次型，它也是相合规范形的一种. 而对于
\[g(X) = x_1^2 + x_2^2 + \cdots + x_{n - 1}^2 + 0 \cdot x_n^2,\]
我们便可以构造出$X_n=(0, 0, \ldots, 0, 1)$使得$g(X)=0$，但去除掉$X_n$生成的子空间和零向量后，剩下的$X$都满足$g(X) > 0$，这是一个半正定二次型. 所以我们可以感觉到正定二次型的条件应当较为严苛，半正定二次型的条件略有放松，但都应该和惯性指数有着密切的联系，这也符合于我们的直觉.

\begin{theorem}{}{}
    设$f(X)$是$n$元实二次型，$A$是$f(X)$的相伴矩阵，$r$为$A$的秩，$r_+$为$A$的正惯性指数. 则
    \begin{enumerate}
        \item $f(X)$是正定二次型当且仅当$r_+=n$；
        \item $f(X)$是半正定二次型当且仅当$r_+=r$.
    \end{enumerate}
\end{theorem}

这一定理证明是显然的，故略去不在此展开. 由此而得到的一个显然的推论如下：

\begin{corollary}{}{正定矩阵的合同标准形}
    $n$阶实对称矩阵$A$是正定矩阵当且仅当其合同于单位阵$E_n$，是半正定矩阵当且仅当其合同于对角阵$\begin{pmatrix}
        E_r & O \\ O & O
    \end{pmatrix}$.
\end{corollary}

考虑正定矩阵$A$，其经过合同变换得到单位阵$E_n$，即$C^\mathrm{T}AC = E_n$，两侧取行列式有$\det(C^\mathrm{T})\det(A)\det(C)=\det(C)^2\det(A)=\det(E_n)=1$，所以$\det(A)>0$，即正定矩阵的行列式大于$0$. 而考虑对正定矩阵做限制，即只取前$r$行$r$列，读者不难证明得到的$r$阶子矩阵也必然是正定矩阵（下面也会给出这一证明），回忆\hyperref[def:子式、主子式、顺序主子式]{顺序主子式的定义}，这启示我们可以通过顺序主子式的符号来判断矩阵的正定性，由此而得到下面的 Sylvester 判别法：

\begin{theorem}{Sylvester 判别法}{Sylvester判别法}
    $n$阶实对称矩阵$A$是正定阵的充分必要条件是$A$的所有顺序主子式的行列式都大于$0$.
\end{theorem}

\begin{proof}
    先证必要性. 设 $A$ 是正定矩阵，记其前 $k$ 行前 $k$ 列构成的顺序主子阵为 $A_k$. 对任意非零向量 $Y=(y_1,\ldots,y_k)^\mathrm{T}$，令 $X=(y_1,\ldots,y_k,0,\ldots,0)^\mathrm{T}$，则 $Y^\mathrm{T}A_kY=X^\mathrm{T}AX>0$. 因此 $A_k$ 也是正定矩阵. 由前面关于正定矩阵的讨论可知 $\det A_k>0$. 所以 $A$ 的所有顺序主子式都大于 $0$.

    再证充分性. 对 $n$ 作归纳. 当 $n = 1$ 时结论显然成立. 下面设结论对 $n - 1$ 阶情形成立，考虑 $n$ 阶实对称矩阵
    \[
        A=\begin{pmatrix}
            a_{11} & \alpha^\mathrm{T} \\
            \alpha & B
        \end{pmatrix},
    \]
    其中 $\alpha\in\mathbf{R}^{n-1}$，$B$ 是 $n - 1$ 阶实对称矩阵. 由第一个顺序主子式大于 $0$ 知 $a_{11}>0$. 定义
    \[
        S=B-\frac{1}{a_{11}}\alpha\alpha^\mathrm{T}.
    \]
    这是一个 $n - 1$ 阶实对称矩阵. 对任意 $k=1,2,\ldots,n-1$，设 $S_k$ 为 $S$ 的前 $k$ 阶顺序主子阵，$A_{k+1}$ 为 $A$ 的前 $k+1$ 阶顺序主子阵，则
    \[
        A_{k+1}=
        \begin{pmatrix}
            a_{11} & \alpha_k^\mathrm{T} \\
            \alpha_k & B_k
        \end{pmatrix},
    \]
    其中 $\alpha_k$ 和 $B_k$ 分别是 $\alpha$ 与 $B$ 的前 $k$ 个分量、前 $k$ 阶顺序主子阵. 对 $A_{k+1}$ 做初等行列变换可得
    \[
        \det A_{k+1}=a_{11}\det S_k.
    \]
    由于 $\det A_{k+1}>0$ 且 $a_{11}>0$，故 $\det S_k>0$. 因而 $S$ 的所有顺序主子式都大于 $0$. 由归纳假设，$S$ 正定.

    现对任意非零向量 $X=(x_1,Y^\mathrm{T})^\mathrm{T}\in\mathbf{R}^n$，有
    \begin{align*}
        X^\mathrm{T}AX
         & = a_{11}x_1^2+2x_1\alpha^\mathrm{T}Y+Y^\mathrm{T}BY \\
         & = a_{11}\left(x_1+\frac{\alpha^\mathrm{T}Y}{a_{11}}\right)^2+Y^\mathrm{T}\left(B-\frac{1}{a_{11}}\alpha\alpha^\mathrm{T}\right)Y \\
         & = a_{11}\left(x_1+\frac{\alpha^\mathrm{T}Y}{a_{11}}\right)^2+Y^\mathrm{T}SY.
    \end{align*}
    若 $Y\neq 0$，则因 $S$ 正定知 $Y^\mathrm{T}SY>0$；若 $Y=0$，则 $x_1\neq 0$，于是第一项大于 $0$. 因而对任意非零 $X$ 都有 $X^\mathrm{T}AX>0$，即 $A$ 正定.
\end{proof}

\begin{example}{}{}
    求$t$的取值范围，使得下面的二次型是正定的：
    \[f(x_1,x_2,x_3,x_4) = x_1^2+4x_2^2+4x_3^2+3x_4^2+2tx_1x_2-2x_1x_3+4x_1x_2.\]
\end{example}

\begin{solution}
    该二次型的相伴矩阵为
    \[
        A=\begin{pmatrix}
            1 & t+2 & -1 & 0 \\
            t+2 & 4 & 0 & 0 \\
            -1 & 0 & 4 & 0 \\
            0 & 0 & 0 & 3
        \end{pmatrix}.
    \]
    由 Sylvester 判别法，只需计算顺序主子式：
    \[
        D_1=1>0,
    \]
    \[
        D_2=\begin{vmatrix}
            1 & t+2 \\
            t+2 & 4
        \end{vmatrix}
        =4-(t+2)^2,
    \]
    \[
        D_3 = \begin{vmatrix}
                 1 & t+2 & -1 \\
                 t+2 & 4 & 0 \\
                 -1 & 0 & 4
             \end{vmatrix} = 4\bigl(3-(t+2)^2\bigr),
    \]
    \[
        D_4=3D_3=12\bigl(3-(t+2)^2\bigr).
    \]
    因而正定的充要条件是 $D_2>0,\qquad D_3>0,\qquad D_4>0$. 其中 $D_3>0$ 已经推出 $D_2>0$ 和 $D_4>0$，故最终只需 $3-(t+2)^2>0$. 即 $-2-\sqrt{3}<t<-2+\sqrt{3}$.
\end{solution}

而通过一些简单的行列对换，我们可以得到这一定理的推论：
\begin{corollary}{}{}
    若 $A$ 是正定矩阵，则：
    \begin{enumerate}
        \item $A$的任一$k$阶主子阵都是正定矩阵；
        \item $A$的所有主子式均为正，特别地，$A$的主对角元素全大于零.
    \end{enumerate}
\end{corollary}

再考虑对正定二次型进行正交变换，得到$f(X) = \lambda_1 x_1^2 + \lambda_2 x_2^2 + \cdots + \lambda_n x_n^2$，其中$\lambda_i$均为$A$的特征值，显然根据上面的推论我们也可以得到如下定理：

\begin{theorem}{}{}
    $n$元实二次型$f(X)$的特征值都是正数当且仅当$f(X)$是正定二次型.
\end{theorem}

不过也可以根据正定矩阵的分解得到以上的结论，对于正定矩阵，我们有如下的引理：

\begin{lemma}{}{}
    设$A$是$n$阶实对称矩阵，那么$A$是正定矩阵的充要条件是存在$n$阶可逆矩阵$C$使得$A=C^\mathrm{T}C$.
\end{lemma}

\begin{proof}
    若存在可逆矩阵 $C$ 使得 $A=C^\mathrm{T}C$，则对任意非零向量 $X$ 有
    \[
        X^\mathrm{T}AX=X^\mathrm{T}C^\mathrm{T}CX=(CX)^\mathrm{T}(CX)=\|CX\|^2>0,
    \]
    因为 $C$ 可逆，故 $CX\neq 0$. 因而 $A$ 正定.

    反之，若 $A$ 正定，则由\autoref{cor:正定矩阵的合同标准形}可知 $A$ 合同于单位阵 $E_n$，即存在可逆矩阵 $P$ 使得 $P^\mathrm{T}AP=E_n$. 两边同乘以 $P^{-\mathrm{T}}$ 与 $P^{-1}$，得到 $A=P^{-\mathrm{T}}P^{-1}$. 令 $C=P^{-1}$，则 $A=C^\mathrm{T}C$.
\end{proof}

设$\lambda$为$A$的任一特征值，$X$为对应的特征向量，则有$AX = \lambda X$，代入引理得到$C^\mathrm{T}CX = \lambda X$，两侧同时左乘$X^\mathrm{T}$，得到
\[ X^\mathrm{T}C^\mathrm{T}CX = (CX)^\mathrm{T}(CX) = \lambda X^\mathrm{T}X, \]
取欧几里得内积，得到 $\|CX\|^2 = \lambda \|X\|^2$. 因为$C$可逆且$X \neq 0$，所以$\lambda > 0$，即$A$的特征值均为正数.

在先前我们介绍过$LU$分解，对于正定矩阵，我们有更加高效的分解方法，即 Cholesky 分解.

\begin{theorem}{}{Cholesky 分解}
    对于正定矩阵$A$，存在唯一的下三角矩阵$L$使得$A=LL^\mathrm{T}$，并且$L$的对角元素均为正数.
\end{theorem}

\begin{proof}
    对 $n$ 作归纳. 当 $n=1$ 时，$A=(a_{11})$ 且 $a_{11}>0$，取 $L=(\sqrt{a_{11}})$ 即可. 下面设结论对 $n-1$ 阶正定矩阵成立. 设
    \[
        A=\begin{pmatrix}
            a_{11} & \alpha^\mathrm{T} \\
            \alpha & B
        \end{pmatrix}
    \]
    为 $n$ 阶正定矩阵. 由正定性知 $a_{11}>0$. 令
    \[
        S=B-\frac{1}{a_{11}}\alpha\alpha^\mathrm{T}.
    \]
    在证明\nameref{thm:Sylvester判别法}时我们知道 $S$ 也是正定矩阵. 由归纳假设，存在唯一的下三角矩阵 $L_1$，其对角元均为正，使得 $S=L_1L_1^\mathrm{T}$. 令
    \[
        L=\begin{pmatrix}
            \sqrt{a_{11}} & 0 \\
            \dfrac{\alpha}{\sqrt{a_{11}}} & L_1
        \end{pmatrix},
    \]
    则 $L$ 为下三角矩阵，且其对角元均为正. 直接计算可得
    \[
        LL^\mathrm{T}
        =
        \begin{pmatrix}
            a_{11} & \alpha^\mathrm{T} \\
            \alpha & \dfrac{1}{a_{11}}\alpha\alpha^\mathrm{T}+L_1L_1^\mathrm{T}
        \end{pmatrix}
        =
        \begin{pmatrix}
            a_{11} & \alpha^\mathrm{T} \\
            \alpha & B
        \end{pmatrix}
        =A.
    \]
    于是存在性成立.

    再证唯一性. 设 $A=LL^\mathrm{T}=\widetilde{L}\widetilde{L}^\mathrm{T}$，其中 $L,\widetilde{L}$ 都是对角元为正的下三角矩阵. 比较 $(1,1)$ 元得
    \[
        l_{11}^2=\widetilde{l}_{11}^2=a_{11}.
    \]
    又因二者都为正，故 $l_{11}=\widetilde{l}_{11}$. 再比较第一列的其余元素，可得对于 $i = 2, 3, \ldots, n$，有 $l_{i1}l_{11}=\widetilde{l}_{i1}\widetilde{l}_{11}$，从而 $l_{i1}=\widetilde{l}_{i1}$. 去掉第一行第一列后，对右下角 $(n-1)$ 阶块继续同样讨论，递归下去便得到 $L=\widetilde{L}$. 唯一性得证.
\end{proof}

半正定矩阵关于以上的结论也是类似的，但相应地有条件的减弱，以及部分等价条件不再成立.
\begin{theorem}{}{}
    若$A$为$n$阶实对称矩阵，则
    \begin{enumerate}
        \item $A$为半正定矩阵等价于$A$的所有特征值非负；
        \item 若$A$为半正定矩阵，则存在矩阵$C$使得$A=C^{\mathrm{T}}C$（注意$C$未必可逆）；
        \item $A$ 为半正定矩阵等价于 $A$ 的各阶主子式非负；
        \item 若 $A$ 为半正定矩阵，则 $A$ 的顺序主子式非负；
        \item 若$A$为半正定矩阵，则存在对角线上元素非负的下三角矩阵$L$使得$A=LL^{\mathrm{T}}$.
    \end{enumerate}
\end{theorem}

根据定理陈述不难看出，$A$ 的顺序主子式非负是半正定矩阵的必要条件，但并非充分条件. 例如
\[A=\begin{pmatrix}
    0 & 0 \\
    0 & -1
\end{pmatrix}\]
的顺序主子式非负，但 $A$ 不是半正定矩阵. 此时，需要将条件加强为各阶主子式非负才是半正定矩阵的充分必要条件. 此外，Cholesky 分解对的唯一性也不再成立，例如
\[A=\begin{pmatrix}
    1 & 0 \\
    0 & 0
\end{pmatrix} = \begin{pmatrix}
    1 & 0 \\
    0 & 0
\end{pmatrix}\begin{pmatrix}
    1 & 0 \\
    0 & 0
\end{pmatrix} = \begin{pmatrix}
    1 & 0 \\
    0 & 1
\end{pmatrix}\begin{pmatrix}
    1 & 0 \\
    0 & 0
\end{pmatrix}\]
都是 $A$ 的 Cholesky 分解. 当然，其余的结论都与正定的情况类似. 下面我们来证明这些结论.

\begin{proof}
    第 1 条显然，故略去证明. 第 2 条可由第 1 条立即得到：若 $A$ 半正定，则 $A=P^\mathrm{T}\diag(\lambda_1,\ldots,\lambda_n)P$，其中 $\lambda_i\geqslant 0$. 令
    \[
        C=\diag(\sqrt{\lambda_1},\ldots,\sqrt{\lambda_n})P,
    \]
    则 $C^\mathrm{T}C=P^\mathrm{T}\diag(\lambda_1,\ldots,\lambda_n)P=A$.

    反过来，若 $A=C^\mathrm{T}C$，则对任意 $X$ 都有
    \[
        X^\mathrm{T}AX=X^\mathrm{T}C^\mathrm{T}CX=(CX)^\mathrm{T}(CX)=\|CX\|^2\geqslant 0,
    \]
    故 $A$ 半正定.

    再证第 3 条. 若 $A$ 半正定，则任取一个主子阵 $B$. 不妨设 $B$ 对应于坐标子空间 $V\subseteq \mathbf{R}^n$. 对任意 $Y\in\mathbf{R}^m$，将它按对应位置嵌入为 $\mathbf{R}^n$ 中的向量 $X$（即在 $V$ 中与 $Y$ 相同，其余位置全补 $0$），则 $Y^\mathrm{T}BY=X^\mathrm{T}AX\geqslant 0$. 所以 $B$ 也是半正定矩阵. 由于 $B$ 的特征值都非负，故 $\det B\geqslant 0$，因而 $A$ 的各阶主子式都非负.

    反之，设 $A$ 的各阶主子式都非负. 对 $n$ 作归纳. 当 $n = 1$ 时结论显然成立. 假设结论对 $n - 1$ 阶情形成立. 考虑 $n$ 阶实对称矩阵 $A=(a_{ij})$. 若所有对角元都为 $0$，则对任意 $i\neq j$，由 $2$ 阶主子式非负可得
    \[
        \begin{vmatrix}
            0 & a_{ij} \\
            a_{ij} & 0
        \end{vmatrix}
        =-a_{ij}^2\geqslant 0,
    \]
    故 $a_{ij}=0$. 于是 $A=O$，自然半正定. 否则若存在某个 $a_{ii} > 0$. 由于同时交换第 $1$ 行与第 $i$ 行、再交换第 $1$ 列与第 $i$ 列不改变各阶主子式都非负这一条件，所以不妨设 $a_{11} > 0$. 写成分块形式
    \[
        A=\begin{pmatrix}
            a_{11} & \alpha^\mathrm{T} \\
            \alpha & B
        \end{pmatrix},
        \qquad
        S=B-\frac{1}{a_{11}}\alpha\alpha^\mathrm{T}.
    \]
    对任意指标集 $I\subseteq\{2,\ldots,n\}$，记 $S_I$ 为 $S$ 的对应主子阵，$A_I$ 为 $A$ 中由第 $1$ 行、第 $1$ 列以及指标集 $I$ 对应的那些行列组成的主子阵，则根据\nameref{thm:Sylvester判别法}的证明过程，有
    \[
        \det S_I=\frac{1}{a_{11}}\det A_I\geqslant 0.
    \]
    因而 $S$ 的各阶主子式都非负. 由归纳假设，$S$ 半正定. 于是对任意 $X=(x_1,Y^\mathrm{T})^\mathrm{T}$，有
    \begin{align*}
        X^\mathrm{T}AX
         & = a_{11}x_1^2+2x_1\alpha^\mathrm{T}Y+Y^\mathrm{T}BY \\
         & = a_{11}\left(x_1+\frac{\alpha^\mathrm{T}Y}{a_{11}}\right)^2+Y^\mathrm{T}SY \geqslant 0.
    \end{align*}
    所以 $A$ 半正定. 第 2 条得证.

    第 4 条是第 3 条的直接推论，因为顺序主子式只是主子式中的一部分.

    最后证第 5 条. 仍对 $n$ 作归纳. 当 $n=1$ 时显然成立. 设结论对 $n-1$ 阶情形成立. 若 $a_{11}=0$，则由半正定性知对任意 $j\geqslant 2$，
    \[
        0\leqslant (e_1+te_j)^\mathrm{T}A(e_1+te_j)=2a_{1j}t+a_{jj}t^2, \qquad \forall t\in\mathbf{R}.
    \]
    将 $t$ 取为充分小的正数与负数，便知 $a_{1j}=0$. 因而第一行与第一列全为 $0$，于是
    \[
        A=\begin{pmatrix}
            0 & 0 \\
            0 & B
        \end{pmatrix},
    \]
    其中 $B$ 半正定. 由归纳假设，存在对角元非负的下三角矩阵 $L_1$ 使得 $B=L_1L_1^\mathrm{T}$. 于是
    \[
        L=\begin{pmatrix}
            0 & 0 \\
            0 & L_1
        \end{pmatrix}
    \]
    满足 $A=LL^\mathrm{T}$.

    若 $a_{11}>0$，则仍令
    \[
        S=B-\frac{1}{a_{11}}\alpha\alpha^\mathrm{T}.
    \]
    对任意 $Y$，取 $x_1=-\dfrac{\alpha^\mathrm{T}Y}{a_{11}}$，由 $A$ 半正定可得
    \[
        0\leqslant \begin{pmatrix}
            -\dfrac{\alpha^\mathrm{T}Y}{a_{11}} \\
            Y
        \end{pmatrix}^\mathrm{T}
        A
        \begin{pmatrix}
            -\dfrac{\alpha^\mathrm{T}Y}{a_{11}} \\
            Y
        \end{pmatrix}
        =Y^\mathrm{T}SY.
    \]
    故 $S$ 半正定. 由归纳假设，存在对角元非负的下三角矩阵 $L_1$ 使得 $S=L_1L_1^\mathrm{T}$. 令
    \[
        L=\begin{pmatrix}
            \sqrt{a_{11}} & 0 \\
            \dfrac{\alpha}{\sqrt{a_{11}}} & L_1
        \end{pmatrix},
    \]
    则 $L$ 的对角元都非负，并且
    \[
        LL^\mathrm{T}
        =
        \begin{pmatrix}
            a_{11} & \alpha^\mathrm{T} \\
            \alpha & \dfrac{1}{a_{11}}\alpha\alpha^\mathrm{T}+L_1L_1^\mathrm{T}
        \end{pmatrix}
        =A.
    \]
    第 5 条得证.
\end{proof}

而到如今，我们回望实内积空间时，会发现所谓的 Gram 矩阵其实就是一个半正定矩阵，内积也可以写作一个二次型的形式. 使用任意一个半正定矩阵都可以定义一个内积，而内积空间的基本性质也可以通过半正定矩阵的性质来推导.但对于复内积空间来说，这就依赖于 Hermite 型和 Hermite 矩阵的性质了，不再过多展开.

\section{标准形的应用}

\subsection{基于相合标准形的分解}

在相抵和相似标准形的讨论中，我们都考虑了基于这些标准形的分解的应用. 本小节继续这个话题，考虑基于相合标准形的分解.

\begin{example}{}{}
    证明：每个秩为$r$的$n\enspace(r<n)$阶实对称矩阵均可表示为$n-r$个秩为$n-1$的实对称矩阵的乘积.
\end{example}

\begin{solution}
    设 $A$ 是秩为 $r$ 的 $n$ 阶实对称矩阵，且 $r<n$. 由\nameref{thm:实谱定理}，存在正交矩阵 $Q$ 使得 $A=Q^\mathrm{T}\Lambda Q$，其中 $\Lambda=\diag(\lambda_1,\ldots,\lambda_r,0,\ldots,0)$ 且 $\lambda_1,\ldots,\lambda_r\neq 0$.

    下面将对角矩阵 $\Lambda$ 分解成 $n-r$ 个秩为 $n-1$ 的对角矩阵的乘积. 定义
    \[
        D_1=\diag(\lambda_1,\ldots,\lambda_r,0,1,\ldots,1),
    \]
    对于 $s=2,\ldots,n-r$，定义
    \[
        D_s=\diag(\underbrace{1,\ldots,1}_{r+s-1},0,1,\ldots,1).
    \]
    每个 $D_s$ 都是实对称矩阵，且恰有一个对角元为 $0$，因此 $r(D_s)=n-1$ 并且显然有 $\Lambda=D_1D_2\cdots D_{n-r}$. 令 $A_s=Q^\mathrm{T}D_sQ$（$s=1,2,\ldots,n-r$），则每个 $A_s$ 都是实对称矩阵，并且由于正交相似保持秩不变，有 $r(A_s)=n-1$. 同时
    \[
        A_1A_2\cdots A_{n-r}
        =Q^\mathrm{T}D_1Q\cdots Q^\mathrm{T}D_{n-r}Q
        =Q^\mathrm{T}(D_1D_2\cdots D_{n-r})Q
        =Q^\mathrm{T}\Lambda Q
        =A.
    \]
    因而 $A$ 可表示为 $n-r$ 个秩为 $n-1$ 的实对称矩阵的乘积.
\end{solution}

\subsection{与特征值有关的进一步结论}

从最优化的角度看，二次型与特征值之间的联系尤为重要. 本节将介绍几条最常用的结论. 在本小节中，若 $A$ 是 $n$ 阶实对称矩阵，我们总把它的特征值按
\[
    \lambda_1\geqslant \lambda_2\geqslant \cdots \geqslant \lambda_n
\]
的顺序排列，并对任意非零向量 $X$ 定义 $A$ 的 Rayleigh 商为
\[
    R_A(X)=\frac{X^\mathrm{T}AX}{X^\mathrm{T}X}.
\]

\begin{theorem}{Rayleigh-Ritz 定理}{Rayleigh-Ritz}
    设 $A$ 是 $n$ 阶实对称矩阵，则对任意非零向量 $X$ 都有
    \[
        \lambda_n\leqslant R_A(X)\leqslant \lambda_1,
    \]
    从而
    \[
        \lambda_n\|X\|^2\leqslant X^\mathrm{T}AX\leqslant \lambda_1\|X\|^2.
    \]
    特别地，
    \[
        \lambda_1=\max_{X\neq 0}R_A(X), \qquad \lambda_n=\min_{X\neq 0}R_A(X).
    \]
\end{theorem}

\begin{proof}
    由\nameref{thm:实谱定理}，存在正交矩阵 $P$ 使得 $P^\mathrm{T}AP=\diag(\lambda_1,\ldots,\lambda_n)$. 对任意非零向量 $X$，令 $X=PY$，则 $Y\neq 0$ 且
    \[
        R_A(X)=\frac{Y^\mathrm{T}(P^\mathrm{T}AP)Y}{Y^\mathrm{T}Y}
        =\frac{\lambda_1y_1^2+\cdots+\lambda_ny_n^2}{y_1^2+\cdots+y_n^2}.
    \]
    右端是 $\lambda_1,\ldots,\lambda_n$ 的加权平均，其中权都非负且和为 $1$，故 $\lambda_n\leqslant R_A(X)\leqslant \lambda_1$. 两边同时乘以 $\|X\|^2=\|Y\|^2$ 便得
    \[
        \lambda_n\|X\|^2\leqslant X^\mathrm{T}AX\leqslant \lambda_1\|X\|^2.
    \]
    当 $X$ 取 $A$ 对应于 $\lambda_1$ 或 $\lambda_n$ 的特征向量时，上下界分别取到，因此
    \[
        \lambda_1=\max_{X\neq 0}R_A(X), \qquad \lambda_n=\min_{X\neq 0}R_A(X).
    \]
\end{proof}

\begin{example}{}{}
    求二次型
    \[
        f(x_1,x_2)=2x_1^2+5x_2^2+4x_1x_2
    \]
    在约束 $x_1^2+x_2^2=1$ 下的最大值与最小值.
\end{example}

\begin{solution}
    该二次型的相伴矩阵为
    \[
        A=\begin{pmatrix}
            2 & 2 \\
            2 & 5
        \end{pmatrix}.
    \]
    在\autoref{ex:正交变换法例1}已经算出它的特征值为 $1,6$. 由于约束条件 $x_1^2+x_2^2=1$ 等价于 $\|X\|=1$，故 $f(X)=X^\mathrm{T}AX=R_A(X)$. 由\autoref{thm:Rayleigh-Ritz}定理立刻得到
    \[
        \min_{\|X\|=1}f(X)=1,\qquad \max_{\|X\|=1}f(X)=6.
    \]
\end{solution}

\begin{theorem}{Courant-Fischer 极小极大定理}{Courant-Fischer}
    设 $A$ 是 $n$ 阶实对称矩阵，则对每个 $k = 1, 2, \ldots, n$，都有
    \[
        \lambda_k = \max_{\dim U=k}\ \min_{0\neq X\in U}R_A(X) = \min_{\dim W=n-k+1}\ \max_{0\neq X\in W}R_A(X).
    \]
\end{theorem}

\begin{proof}
    由\nameref{thm:实谱定理}，存在 $A$ 的一组标准正交特征向量 $\xi_1,\ldots,\xi_n$，满足 $A\xi_i=\lambda_i\xi_i,\qquad i=1,\ldots,n$. 考虑任意 $k$ 维子空间 $U$（$n - k + 1$ 维子空间 $U$ 的情况留作习题），因为
    \[
        \dim U+\dim\spa(\xi_k,\xi_{k+1},\ldots,\xi_n)=k+(n-k+1)=n+1,
    \]
    所以 $U\cap \spa(\xi_k,\xi_{k+1},\ldots,\xi_n)\neq \{\vec{0}\}$. 取其中非零向量 $X=a_k\xi_k+\cdots+a_n\xi_n$，则
    \[
        R_A(X)=\frac{\lambda_ka_k^2+\cdots+\lambda_na_n^2}{a_k^2+\cdots+a_n^2}\leqslant \lambda_k.
    \]
    因而 $\min_{0\neq X\in U}R_A(X)\leqslant \lambda_k$. 对所有 $k$ 维子空间取最大，得到
    \[
        \max_{\dim U=k}\ \min_{0\neq X\in U}R_A(X)\leqslant \lambda_k.
    \]
    另一方面，取 $U_0=\spa(\xi_1,\ldots,\xi_k)$，则对任意非零向量 $X=a_1\xi_1+\cdots+a_k\xi_k\in U_0$ 都有
    \[
        R_A(X)=\frac{\lambda_1a_1^2+\cdots+\lambda_ka_k^2}{a_1^2+\cdots+a_k^2}\geqslant \lambda_k.
    \]
    故
    \[
        \min_{0\neq X\in U_0}R_A(X)\geqslant \lambda_k.
    \]
    再取 $X=\xi_k$，便知这里的最小值恰为 $\lambda_k$. 因而
    \[
        \max_{\dim U=k}\ \min_{0\neq X\in U}R_A(X)=\lambda_k.
    \]
\end{proof}

下面的\nameref{thm:Cauchy交错}和\nameref{thm:Weyl}是 Courant-Fischer 极小极大定理的应用.

\begin{theorem}{Cauchy 交错定理}{Cauchy交错}
    设 $A$ 是 $n$ 阶实对称矩阵，特征值为 $\lambda_1\geqslant \lambda_2\geqslant \cdots \geqslant \lambda_n$. 设 $B$ 是 $A$ 的一个 $m$ 阶主子阵，特征值为 $\mu_1\geqslant \mu_2\geqslant \cdots \geqslant \mu_m$. 则对每个 $i=1,2,\ldots,m$，都有
    \[
        \lambda_i\geqslant \mu_i\geqslant \lambda_{n-m+i}.
    \]
\end{theorem}

不难看出，Cauchy 交错定理表明：若 $A$ 正定，则其任一主子阵也正定.

\begin{proof}
    不妨设 $B$ 对应于坐标子空间 $V\subseteq \mathbf{R}^n$. 对任意 $Y\in\mathbf{R}^m$，将它按对应位置嵌入为 $\mathbf{R}^n$ 中的向量 $\widetilde{Y}$（即在 $V$ 中与 $Y$ 相同，其余位置全补 $0$），则
    \[
        Y^\mathrm{T}BY=\widetilde{Y}^\mathrm{T}A\widetilde{Y}, \qquad Y^\mathrm{T}Y=\widetilde{Y}^\mathrm{T}\widetilde{Y}.
    \]
    因而 $B$ 在 $Y$ 上的 Rayleigh 商恰等于 $A$ 在 $\widetilde{Y}$ 上的 Rayleigh 商. 由\nameref{thm:Courant-Fischer}，
    \[
        \mu_i=\max_{\substack{\dim U=i \\ U\subseteq V}}\ \min_{0\neq X\in U}R_A(X).
    \]
    由于这里取最大时只在 $V$ 的子空间中取，而 $A$ 的对应极大值是在所有 $i$ 维子空间中取，所以
    \[
        \mu_i\leqslant \max_{\dim U=i}\ \min_{0\neq X\in U}R_A(X)=\lambda_i.
    \]

    同样地，由\nameref{thm:Courant-Fischer}，
    \[
        \mu_i=\min_{\substack{\dim W=m-i+1 \\ W\subseteq V}}\ \max_{0\neq X\in W}R_A(X).
    \]
    而
    \[
        \lambda_{n-m+i}=\min_{\dim W=m-i+1}\ \max_{0\neq X\in W}R_A(X).
    \]
    由于后者是在所有维数为 $m-i+1$ 的子空间上取最小，前者只是在其中一部分子空间上取最小，所以 $\lambda_{n-m+i}\leqslant \mu_i$. 综合上述讨论即得 $\lambda_i\geqslant \mu_i\geqslant \lambda_{n-m+i}$.
\end{proof}

\begin{theorem}{Weyl 不等式}{Weyl}
    设 $A,B$ 是 $n$ 阶实对称矩阵，它们的特征值分别按降序排列为 $\alpha_1\geqslant \cdots \geqslant \alpha_n, \beta_1\geqslant \cdots \geqslant \beta_n$. 设 $A+B$ 的特征值按降序排列为 $\gamma_1\geqslant \cdots \geqslant \gamma_n$. 则
    \begin{enumerate}
        \item 若 $i+j-1\leqslant n$，则 $\gamma_{i+j-1}\leqslant \alpha_i+\beta_j$;
        \item 若 $i+j-n\geqslant 1$，则 $\gamma_{i+j-n}\geqslant \alpha_i+\beta_j$.
    \end{enumerate}
\end{theorem}

\begin{proof}
    设 $\xi_1,\ldots,\xi_n$ 是 $A$ 的一组标准正交特征向量，满足 $A\xi_s=\alpha_s\xi_s$；设 $\eta_1,\ldots,\eta_n$ 是 $B$ 的一组标准正交特征向量，满足 $B\eta_t=\beta_t\eta_t$.

    下面证明第 1 条，第 2 条思路类似，故留作习题. 令 $U=\spa(\xi_i,\xi_{i+1},\ldots,\xi_n), V=\spa(\eta_j,\eta_{j+1},\ldots,\eta_n)$. 则 $\dim U=n-i+1, \dim V=n-j+1$，从而
    \[ \dim(U\cap V)\geqslant (n-i+1)+(n-j+1)-n=n-i-j+2. \]
    当 $i+j-1\leqslant n$ 时，上式右端正好是 $n-(i+j-1)+1$. 取 $U\cap V$ 的一个维数为 $n-(i+j-1)+1$ 的子空间 $W$. 对任意非零 $X\in W$，根据 $X\in U$ 与 $X\in V$，分别有 $R_A(X)\leqslant \alpha_i, R_B(X)\leqslant \beta_j$，故
    \[
        R_{A+B}(X)=R_A(X)+R_B(X)\leqslant \alpha_i+\beta_j.
    \]
    因而
    \[
        \max_{0\neq X\in W}R_{A+B}(X)\leqslant \alpha_i+\beta_j.
    \]
    由\nameref{thm:Courant-Fischer}，
    \[
        \gamma_{i+j-1}=\min_{\dim W'=n-(i+j-1)+1}\ \max_{0\neq X\in W'}R_{A+B}(X)\leqslant \alpha_i+\beta_j.
    \]
\end{proof}

\begin{corollary}{}{Weyl 不等式的推论1}
    设 $A,B$ 是 $n$ 阶实对称矩阵，特征值分别按降序排列为 $\alpha_1\geqslant \cdots \geqslant \alpha_n$ 与 $\beta_1\geqslant \cdots \geqslant \beta_n$，设 $A+B$ 的特征值为 $\gamma_1\geqslant \cdots \geqslant \gamma_n$. 则对每个 $k=1,2,\ldots,n$，都有
    \[
        \alpha_k+\beta_n\leqslant \gamma_k\leqslant \alpha_k+\beta_1.
    \]
\end{corollary}

\begin{proof}
    在\nameref{thm:Weyl}中分别取 $(i,j)=(k,1)$ 与 $(i,j)=(k,n)$ 即可.
\end{proof}

\begin{corollary}{}{Weyl 不等式的推论2}
    设 $A,B$ 是 $n$ 阶实对称矩阵. 若 $B-A$ 半正定，则对每个 $k=1,2,\ldots,n$，都有
    \[
        \lambda_k(A)\leqslant \lambda_k(B).
    \]
\end{corollary}

\begin{proof}
    记 $C=B-A$. 由假设知 $C$ 半正定，所以它的最小特征值满足 $\lambda_n(C)\geqslant 0$. 又 $B=A+C$，故由\autoref{cor:Weyl 不等式的推论1}可得 $\lambda_k(B)\geqslant \lambda_k(A)+\lambda_n(C)\geqslant \lambda_k(A)$.
\end{proof}

\begin{summary}
    本讲围绕实数域上二次型展开，此时二次型与对称矩阵一一对应. 首先我们讨论了如何求二次型的标准形，本质上就是把相伴矩阵作合同对角化. 这可以通过初等变换、配方法或正交变换来完成. 在实数域上，合同对角化进一步导向相合规范形与惯性定理：实对称矩阵的正、负、零惯性指数是合同不变量. 由此可以把正定、半正定等保号性质完全转化为规范形、特征值与矩阵分解的条件，并得到 Sylvester 判别法、特征值判别法以及 Cholesky 分解等一系列常用工具. 在应用部分，一方面可以借助合同标准形处理矩阵分解问题；另一方面，Rayleigh-Ritz、Courant-Fischer 极小极大定理、Cauchy 交错定理与 Weyl 不等式把二次型和特征值估计紧密联系起来，这正是后续微积分、最优化和数值分析中最常见的线性代数基础.
\end{summary}

\begin{exercise}
    % \exquote[]{}

    \begin{exgroup}
        \item 设$A=\begin{pmatrix}
            1 & 2 & 0 \\ 2 & 1 & 0 \\ 0 & 0 & 3
        \end{pmatrix},\enspace B=\begin{pmatrix}
            -2 & 0 & 0 \\ 0 & 2 & 1 \\ 0 & 1 & 2
        \end{pmatrix}$，判断$A$与$B$是否相合.

        \item 回答如下有关合同性质的问题：
        \begin{enumerate}
            \item 证明：合同是等价关系；
            \item 证明：在复数域上，两个矩阵相合当且仅当它们相抵；并基于此举例说明：在复数域上相合的矩阵不一定在实数域上相合；
            \item 举例说明：$A \simeq B$ 不能得到 $A^m \simeq B^m$；
            \item 证明：当 $A,B$ 是实对称矩阵时，$A\simeq B$ 当且仅当 $A^m\simeq B^m$；
            \item 证明：若可逆矩阵 $A, B$ 满足 $A\simeq B$，则 $A^{-1}\simeq B^{-1}$；
            \item 证明：若 $A_1\simeq A_2$ 且 $B_1\simeq B_2$，则
                  \[
                      \begin{pmatrix}
                          A_1 & O \\ O & B_1
                      \end{pmatrix}\simeq\begin{pmatrix}
                          A_2 & O \\ O & B_2
                      \end{pmatrix};
                  \]
            \item 设$A\simeq B$，$C\simeq D$，且它们都是$n$阶实对称矩阵，判断 $A+C\simeq B+D$ 是否成立，正确请说明理由，不正确请举出反例；
            \item  判断：矩阵相似是否一定合同？矩阵合同是否一定相似？对于实对称矩阵上述论断又是否正确呢？正确请说明理由，不正确请举出反例.
        \end{enumerate}
    \end{exgroup}

    \begin{exgroup}
        \item 证明：$A\simeq B$表明$A$可以每次做相同的初等行列变换得到$B$，反之亦然

        \item 不直接求偏导，利用 Rayleigh-Ritz 定理求二次型 $2x_1^2+5x_2^2+6x_1x_2$ 在单位圆 $x_1^2+x_2^2=1$ 上的最大值与最小值.

        \item 完善\nameref{thm:Courant-Fischer}的证明.

        \item 完善\nameref{thm:Weyl}的证明.
    \end{exgroup}

    \begin{exgroup}
        \item 设$n$元二次型$f(x_1,x_2,\ldots,x_n)=l_1^2+\cdots+l_p^2-l_{p+1}^2-\cdots-l_{p+q}^2$，其中$l_i\enspace (i=1,2,\ldots,p+q)$是关于$x_1,x_2,\ldots,x_n$的一次齐次式. 证明：$f(x_1,x_2,\ldots,x_n)$的正惯性指数$\leqslant p$，负惯性指数$\leqslant q$；

        \item 已知$A$为$m$阶实对称矩阵，$C$为$m\times n$实矩阵，证明：$C^\mathrm{T}AC$的正负惯性指数分别小于等于$A$的正负惯性指数.

        \item 设 $A$ 是 $n$ 阶实对称矩阵，证明：对任意实数 $t$，$A+tE_n$ 的特征值恰为 $\lambda_1(A)+t,\ldots,\lambda_n(A)+t$. 由此求使得 $A+tE_n$ 半正定的最小实数 $t$.

        \item 设 $A,B$ 都是 $n$ 阶实对称矩阵. 利用 Weyl 不等式证明：
        \[
            |\lambda_k(A)-\lambda_k(B)|\leqslant \|A-B\|_2, \qquad k=1,2,\ldots,n,
        \]
        其中 $\|\,\cdot\,\|_2$ 表示谱范数.

        \item 设 $A,B$ 都是 $n$ 阶实对称矩阵，且 $B-A$ 半正定. 证明：对任意非零向量 $X$，若 $X^\mathrm{T}AX>0$，则 $X^\mathrm{T}BX>0$；并说明这与特征值单调性之间的关系.
    \end{exgroup}
\end{exercise}
